\section{实施风险与应对}
\label{sec:discussion}
\subsection{数据域风险}
计划在数据准备阶段区分自然场景图像与近红外座舱图像，为各域保留独立验证集合。如果候选只改善某一来源，先检查采样比例、归一化与增强，再决定是否采用。数据清单与划分固定后保留版本，避免优化过程中改变评价对象。

\subsection{中性表情混淆}
拟将中性识别作为专门检查项目。若中性召回或中性到悲伤、愤怒的错误数量未达到预期，先核对真实标签与人脸裁剪，再利用 PPB-Emo 中性样本和匿名域训练进行调整。模型输出概率与时序结果分别观察，区分单帧分类和连续决策问题。

\subsection{关键点与速度}
低对比度或遮挡输入可能影响人脸与关键点定位。拟为检测、匿名化和分类设置明确状态，失败时停止当前帧的分类，不以原脸替代匿名输入。性能问题则通过分段计时定位，在保持输入顺序和表情准确性的前提下采用缓存、检测复用与 GPU 推理。

\subsection{验证口径}
计划将每项结果与数据来源、人员划分和运行设置关联保存。身份评价固定自建模型，表情评价固定真实标签，帧率评价固定模型计时区间。材料汇总时按同一记录编写图表，使方法选择、实现与验收能够前后对应。
